\section{Subsquare and congruence obstructions}

The FFF property is hereditary for Latin subsquares in a stronger,
componentwise sense.

\begin{theorem}[Subsquare monotonicity]\label{thm:subsquare-monotonicity}
If $M$ is a Latin subsquare of $L$, then
\[
  \pat(M)\leq \pat(L)
\]
componentwise, with $F<T$.
\end{theorem}

\begin{proof}
Let the rows, columns, and symbols of $M$ be $R_0,C_0,S_0$.  For
$a,b\in R_0$, the ambient row permutation $r_b^{-1}r_a$ preserves $C_0$:
if $c\in C_0$, then $L(a,c)\in S_0$, and row $b$ contains this symbol at a
unique column of $C_0$.  Its restriction is precisely the induced row
permutation of $M$.  Thus every odd row cycle of $M$ is also an odd row
cycle of $L$.  Interchanging rows and columns proves the column statement.

For $u,v\in S_0$ and $c\in C_0$, the row $s_u(c)$ lies in $R_0$.  In that
row, the unique occurrence of $v$ lies in a column $d\in C_0$.  Hence
$s_v^{-1}s_u$ preserves $C_0$ and restricts to the symbol permutation of
$M$.  This proves the symbol statement.
\end{proof}

\begin{corollary}\label{cor:no-odd-subsquare}
Every Latin subsquare of an FFF square is FFF.  In particular, an FFF square
contains no subsquare of odd order greater than one.
\end{corollary}

\begin{proof}
The first assertion is Theorem~\ref{thm:subsquare-monotonicity} with
$\pat(L)=\FFF$.  The second follows from Proposition~\ref{prop:odd-order}.
\end{proof}

Together with Computational Theorem~\ref{cthm:small-orders}, the same implication
excludes order-$6$ subsquares; that additional exclusion inherits the
computational status of the order-$6$ theorem.

There is an exact linear-algebraic description of all half-order
subsquares.  Let $M_L$ be the $3n$-by-$n^2$ binary incidence matrix whose
rows are the row, column, and symbol lines and whose columns are the occupied
cells of $L$.  Let
\[
 U=\{(\alpha\mathbf1_R,\beta\mathbf1_C,\gamma\mathbf1_S):
       \alpha+\beta+\gamma=0\}
 \subseteq\ker M_L^{\mathsf T}
\]
and put
\[
 \delta_2(L)=(3n-2)-\operatorname{rank}_{\mathbb F_2}M_L.
\]

\begin{theorem}[Binary line kernel and half-order subsquares]
\label{thm:binary-kernel-subsquares}
Individual embedded half-order Latin subsquares of $L$ are in bijection with
\[
 \ker M_L^{\mathsf T}\setminus U.
\]
Equivalently, the nonzero classes in
$\ker M_L^{\mathsf T}/U$ parametrize their complementary four-packs.  Thus
their number is exactly
\begin{equation}
 4\bigl(2^{\delta_2(L)}-1\bigr).\label{eq:half-order-count}
\end{equation}
If $n=2m$, where $m>1$ is odd, then any one F-view forces
\begin{equation}
 \operatorname{rank}_{\mathbb F_2}M_L=3n-2.\label{eq:fff-binary-line-rank}
\end{equation}
In particular, \eqref{eq:fff-binary-line-rank} holds for every FFF square of order congruent to $2$
modulo $4$ and greater than $2$.
\end{theorem}

\begin{proof}
A left-kernel vector is a triple $(a,b,c)$ satisfying
\begin{equation}
 a_r+b_j+c_{L(r,j)}=0.\label{eq:binary-line-dependency}
\end{equation}
If one coordinate, say $a$, is constant, then fixing a
column and varying its rows shows that $c$ is constant; fixing a row then
makes $b$ constant.  Hence every vector outside $U$ has all three
coordinates nonconstant.

Fix a row $r$.  Since that row permutes the symbols,
\eqref{eq:binary-line-dependency} identifies the
multiset of $c$-values with that of $a_r+b$.  Choosing rows with $a_r=0$
and $a_r=1$ gives
\[
 \operatorname{wt}(c)=\operatorname{wt}(b)
 =n-\operatorname{wt}(b).
\]
Thus $n=2m$ and $b,c$ each have weight $m$.  The analogous fixed-column
argument gives $\operatorname{wt}(a)=m$.  For their binary fibres
$A_x,B_y,C_z$, equation~\eqref{eq:binary-line-dependency} now gives
\[
 L(A_x,B_y)\subseteq C_{x+y}\qquad(x,y\in\mathbb F_2).
\]
Every set has size $m$; Latinness makes each of the four displayed blocks an
order-$m$ Latin subsquare.

Adding an element of $U$ complements zero or two of the three partitions and
selects a different block in the same four-pack.  Conversely, if
$R_0\times C_0$ is a half-order subsquare with symbol set $S_0$, row and
column counting forces the complementary pattern
\[
 R_0\times C_0,\ R_1\times C_1\longrightarrow S_0,
 \qquad R_0\times C_1,\ R_1\times C_0\longrightarrow S_1.
\]
The three complement indicators satisfy
\eqref{eq:binary-line-dependency}, and the chosen block recovers
them uniquely.  This proves both bijections.  Since $U$ has dimension two
and the quotient has dimension $\delta_2(L)$, counting its nonzero classes
proves the counting formula.

Finally suppose that $m>1$ is odd and $\delta_2(L)>0$.  The preceding
construction gives an order-$m$ subsquare.  Two distinct rows inside it
induce a fixed-point-free permutation preserving its $m$ columns.  Such a
permutation on an odd set has a nontrivial odd cycle, contradicting an F row
view.  The two parastrophes give the same conclusion for either other view,
which proves the rank formula.  The exception $m=1$ is real: the order-$2$ FFF square has
four singleton subsquares and $\delta_2=1$.
\end{proof}

The complete finite audit independently constructs both sides of the
bijection on $9778$ tables and obtains exact set equality in every case.  In
particular, the $230$ order-$8$ FFF representatives have
\[
 \delta_2=0,1,2,3\quad\text{with multiplicities}\quad63,156,10,1,
\]
which explains their half-order-subsquare distribution
$0^{63}4^{156}12^{10}28^1$.  A hypothetical order-$10$ FFF square must have
binary line rank $28$ and no order-$5$ subsquare.  This condition is not
sufficient: all $135$ tracked order-$10$ diagnostic tables already satisfy
it, none is FFF, and that corpus is not exhaustive.

The same kernel equation has a uniform interpretation over every prime
field, now in terms of quotient homotopies rather than individual
subsquares.  Moorhouse's character formula for loop-coordinatized
$3$-nets gives the general identity
$\operatorname{rank}_{p}M_L=3n-2-\dim\operatorname{Hom}(G,\mathbb F_p)$
and identifies its defect with elementary-abelian quotient
characters~\cite[Theorem~4.2]{Moorhouse1991BruckNets}.  We record a direct
three-sorted proof in our notation because the individual cyclic quotient is
the input needed for the FFF corollary below.
The author's publication page flags an erratum for the original article; the
reference entry therefore points to the author-hosted article copy and the
separate erratum listing.  The direct proof given here, rather than an unexamined appeal
to the original pagination, is the logical basis for the specialization used
below.

\begin{theorem}[Prime-field line quotients]
\label{thm:prime-field-line-quotients}
Let $p$ be prime and let
$(a,b,c)\in\ker_{\mathbb F_p}M_L^{\mathsf T}$ be nonconstant modulo the
two-dimensional constant subspace.  Then $p\mid n$, each of $a,b,c$ takes
every value exactly $n/p$ times, and their fibres satisfy
\[
 L(R_x,C_y)\subseteq S_{-x-y}\qquad(x,y\in\mathbb F_p).
\]
After a suitable isotopy these three fibre partitions become one quasigroup
congruence with $p$ blocks of size $n/p$ and quotient isotopic to $C_p$.

Over characteristic zero, $M_L$ has rank $3n-2$.  Every prime dividing a
nonzero Smith invariant factor of the integer matrix $M_L$ divides $n$.
\end{theorem}

\begin{proof}
If one of $a,b,c$ is constant, fixing a line in either other sort and using
Latinness makes the remaining two constant.  Let $\mu_b(t)$ and $\mu_c(t)$
be value multiplicities.  Since every row permutes the symbols,
equation~\eqref{eq:binary-line-dependency}, now read in $\mathbb F_p$, gives
\[
 \mu_c(t)=\mu_b(-a_r-t).
\]
Comparing rows with $a_r-a_{r'}\ne0$ makes $\mu_b$ invariant under a nonzero
translation.  That translation generates the additive group of
$\mathbb F_p$, so $\mu_b$ is uniform.  Thus $p\mid n$ and both $b,c$ are
balanced; the column argument balances $a$.  The displayed fibre rule is the
kernel equation itself.

Choose bijections from one partitioned $n$-set to the row, column, and symbol
sets, matching equally labelled fibres.  Transporting $L$ by these three
bijections gives an isotope whose block operation is $x*y=-x-y$ on
$\mathbb F_p$.  It is Latin and isotopic to addition, so block equivalence is
compatible with multiplication and both divisions and is a quasigroup
congruence.

In characteristic zero a finite nonempty multiset cannot be invariant under
a nonzero translation, because its iterated translates would have infinite
support.  Hence only the two constant dependencies remain.  Finally, if a
prime divides a nonzero Smith invariant factor, reduction modulo that prime
lowers the rank; the first part then forces the prime to divide $n$.
\end{proof}

An eight-worker audit checks this theorem over
$\mathbb F_2,\mathbb F_3,\mathbb F_5,\mathbb F_7,\mathbb F_{11}$ on $9778$
tables.  It verifies all $1160$ nonuniversal quotient classes and a separate
implementation reconstructs $48890$ modular ranks, with no errors.

We next use congruences in the quasigroup algebra with multiplication and the
two division operations.  All congruence blocks have equal size and form a
Latin quotient.

\begin{theorem}[Congruence obstructions]\label{thm:congruence-obstructions}
Let $L$ be a finite Latin square.
\begin{enumerate}[label=(\roman*)]
\item If $L$ has a congruence with odd block size greater than one, then
  $\pat(L)=\TTT$.
\item If the congruence blocks have size two and the quotient is $Q$, then
  $\pat(L)=\pat(Q)$ componentwise.
\end{enumerate}
\end{theorem}

\begin{proof}
For (i), choose distinct rows $x,x'$ in one congruence block $A$ and any
column block $B$.  The two left translations map $B$ bijectively onto the
same product block.  Thus $L_{x'}^{-1}L_x$ preserves $B$ and is
fixed-point-free there.  Since $|B|$ is odd, its restriction has a
nontrivial odd cycle.  Applying the same argument to the two relevant
parastrophes proves the column and symbol bits.

For (ii), label every two-element block by $(i,0),(i,1)$, where $i\in Q$.
The unique form of a binary Latin square gives a function
$f\colon Q^2\to\mathbb F_2$ such that
\[
  (i,a)*(j,b)=(i*j,a+b+f(i,j)).
\]
Rows in the same quotient block induce only transpositions inside the column
fibres.  For rows in different quotient blocks, the lift above a quotient
cycle of length $\ell$ is either two cycles of length $\ell$ or one cycle of
length $2\ell$, according to its binary voltage.  Thus an ambient odd cycle
can occur only above an odd quotient cycle.

Conversely, suppose quotient rows $i,i'$ induce an odd cycle $C$ of length
$\ell>1$.
Above a quotient column $j\in C$, the lifted row permutation comparing
$(i,a)$ with $(i',a')$ has the form
\[
 (j,b)\longmapsto
 \bigl(\pi(j),b+a+a'+f(i,j)+f(i',\pi(j))\bigr).
\]
Its voltage around $C$ is $a+a'+S$, because $\ell$ is odd, where $S$ is
independent of $a,a'$.  Choose $a+a'=S$.  The voltage is then zero, so the
lift of $C$ consists of two cycles of the same odd length $\ell$.  Hence a
quotient row witness lifts to an ambient row witness.  This proves equality
of the row bits.  The same complete argument in the parastrophes proves
equality of the column and symbol bits.
\end{proof}

\begin{corollary}[Order-$2p$ simplicity]\label{cor:2p-simple}
Let $p$ be an odd prime.  Every quasigroup of order $2p$ with at least one
F view is congruence-simple, and every isotope of it is congruence-simple.
\end{corollary}

\begin{proof}
A nontrivial congruence has block-size/quotient-order pair $(p,2)$ or
$(2,p)$.  The first is excluded by
Theorem~\ref{thm:congruence-obstructions}(i).  In the second, the quotient
has pattern TTT by Proposition~\ref{prop:odd-order}, so (ii) excludes it.
Both congruence cases force TTT, contradicting even one F view.
Finally, the full view pattern is isotopy invariant by
Proposition~\ref{prop:invariance}, and the same argument applies to every isotope.
\end{proof}

\begin{corollary}[Order-$2p$ line-incidence saturation]
\label{cor:2p-line-saturation}
Let $p>1$ be an odd prime and let $L$ have order $2p$.  If even one view of
$L$ is F, then $M_L$ has rank $6p-2$ over every field, all its nonzero Smith
invariant factors are $1$, and
\[
 \mathbb Z^{6p}/\operatorname{im}_{\mathbb Z}M_L\cong\mathbb Z^2.
\]
In particular, this holds for every hypothetical order-$10$ FFF square.
\end{corollary}

\begin{proof}
By Theorem~\ref{thm:prime-field-line-quotients}, a torsion prime must divide
$2p$.  A characteristic-two defect yields, after isotopy, a congruence with
two blocks of odd size $p$, so
Theorem~\ref{thm:congruence-obstructions}(i) gives pattern TTT.  A
characteristic-$p$ defect yields blocks of size two; part (ii) identifies the
pattern with that of the odd-order quotient, again TTT by
Proposition~\ref{prop:odd-order}.  Isotopy invariance contradicts the assumed
F view in either case.  Thus no Smith prime remains, proving saturation and
the displayed cokernel.
\end{proof}

Thus a hypothetical order-$10$ FFF square must be isotopically
congruence-simple.  This removes every congruence-imprimitive quasigroup
construction, but it does not rule out simple quasigroups and therefore does
not decide order $10$.

The binary equality iterates through congruence series.

\begin{corollary}[Prime-factor congruence series]
\label{cor:prime-congruence-series}
Suppose that a quasigroup admits a chain of congruence quotients ending in the
one-element quasigroup, with every congruence block size prime.  It is FFF if
all those primes are $2$, and it has pattern TTT if at least one is odd.
\end{corollary}

\begin{proof}
A binary step preserves the entire pattern by
Theorem~\ref{thm:congruence-obstructions}(ii).  At an odd step,
Theorem~\ref{thm:congruence-obstructions}(i) gives TTT independently of the
lower quotient.  Iterating proves both statements.
\end{proof}

This proves the power-of-two conjecture for this congruence-decomposable
class, but leaves congruence-simple quasigroups outside its scope.

The same theorem also gives a census-free nongroup example.

\begin{theorem}[Explicit nongroup FFF loop]
\label{thm:explicit-nongroup-loop}
On $\mathbb Z_4\times\mathbb F_2$, define
\[
 (i,a)*(j,b)=
 \bigl(i+j,\,a+b+\mathbf 1_{(i,j)=(1,1)}\bigr),
\]
with coordinates taken modulo $4$ and $2$, respectively.  This is an FFF
loop of order $8$ that is not group-isotopic.
\end{theorem}

\begin{proof}
The identity is $(0,0)$, and the first-coordinate fibres form a congruence
with quotient $C_4$.  The quotient is FFF by
Corollary~\ref{cor:group-fff}, so binary pattern equality makes the loop
FFF.  It is not associative because
\[
 ((1,0)*(1,0))*(2,0)=(0,1),\qquad
 (1,0)*((1,0)*(2,0))=(0,0).
\]
At the identity base row, the basis family in
Proposition~\ref{prop:closure} is the family of left translations.
Closure would force $L_xL_y=L_{x*y}$ and hence associativity.  The displayed
failure is therefore a closure failure, so the loop is not group-isotopic.
\end{proof}
